\documentclass[11pt,a4paper]{article}
\usepackage[T1]{fontenc}
\usepackage{lmodern,amsmath,amssymb,amsthm,mathtools}
\usepackage[margin=25mm,headheight=15pt]{geometry}
\usepackage{booktabs,array,tabularx}
\usepackage[dvipsnames]{xcolor}
\usepackage{microtype,enumitem,fancyhdr,needspace}
\usepackage[colorlinks=true,linkcolor=MidnightBlue,citecolor=MidnightBlue,urlcolor=MidnightBlue]{hyperref}
\usepackage{bookmark}
\hypersetup{pdftitle={Same-Trajectory Native Credit in the Reciprocal Chiral Exchange Candidate},
pdfauthor={Hilmir Frímann Halldórsson},
pdfsubject={A finite-coupling theorem for an initially empty carrier and receiver. Research manuscript v0.1, draft for manuscript review. 3 November 2026.}}
\setlength{\parskip}{4pt}\setlength{\parindent}{0pt}\setlength{\emergencystretch}{3em}
\setlist{itemsep=2pt,topsep=4pt}\setlist[itemize]{label=$\bullet$}
\pagestyle{fancy}\fancyhf{}
\fancyhead[L]{\small Same-trajectory native credit}\fancyhead[R]{\small v0.1 draft}
\fancyfoot[C]{\thepage}\renewcommand{\headrulewidth}{0.3pt}
\numberwithin{equation}{section}
\newtheorem{theorem}{Theorem}[section]
\newtheorem{lemma}[theorem]{Lemma}
\newcommand{\R}{\mathbb R}\newcommand{\C}{\mathbb C}
\newcommand{\B}{\mathcal B}\newcommand{\LL}{\mathcal L}
\newcommand{\EE}{\mathcal E}\newcommand{\dd}{\,\mathrm d}
\newcommand{\norm}[1]{\left\|#1\right\|}
\newcolumntype{Y}{>{\raggedright\arraybackslash}X}
\newcolumntype{P}[1]{>{\raggedright\arraybackslash}p{#1}}
\begin{document}
\hypersetup{pageanchor=false}
\begin{titlepage}
\vspace*{5mm}
{\small\color{MidnightBlue} TRI-OCTAGON MATHEMATICAL RESEARCH\par}
\vspace{8mm}
{\LARGE\bfseries\raggedright Same-Trajectory Native Credit\\
in the Reciprocal Chiral\\Exchange Candidate\par}
\vspace{6mm}
{\large A finite-coupling theorem for an initially empty carrier and receiver\par}
\vspace{7mm}
{\large Hilmir Fr\'imann Halld\'orsson\par}
\vspace{3mm}
{Research manuscript v0.1\quad 3 November 2026\par}
{\itshape Draft for manuscript review\par}
\vspace{5mm}
\begin{abstract}
The reciprocal chiral exchange candidate for the Tri-Octagon model couples a native state to a periodic carrier and a reversible receiver through two finite distributed profiles. Its earlier analysis established initial supply from an empty carrier and receiver, and native credit from separately prepared states. We prove a sufficient finite-coupling theorem that connects these signs on one unchanged solution. Fix the admissible geometry, positive storage scales and native coupling, and a launch from the specified coherent family with nonzero chiral drive. An explicit finite receiver-coupling threshold ensures initial positive total native power and receiver charging, followed by strictly negative total native-to-carrier power at a specified later time. The proof uses exact time and state scaling, a weighted external-energy estimate, and an oscillator comparison whose propagation and native-drift errors are bounded at finite coupling. Only first weak derivatives of the tent profiles are required. Continuity supplies a later interval of positive native-storage recovery without resetting the state or changing coefficients. Cumulative native debit remains nonnegative, and the coherent-core invariant ceiling persists. The result establishes a behavior of the published candidate, with a dynamically active receiver; it does not classify all positive couplings or establish a substantial recovered fraction. The witness describes fast collocated exchange rather than a completed spatial circuit or calibrated physical energy transport.
\end{abstract}
\vfill
{\small Research author: Hilmir Fr\'imann Halld\'orsson (pzychozen).
GPT formulated and internally reviewed the underlying analytical investigation;
Codex independently verified the argument and prepared this manuscript.
The accepted mathematics and its unchanged constants are the basis of this draft.
Review of the final manuscript by GPT and the author is pending.}
\end{titlepage}
\setcounter{page}{2}\hypersetup{pageanchor=true}
\section{From prepared credit to credit on one solution}\label{sec:introduction}
The reciprocal chiral exchange extension of the Tri-Octagon model provides an explicit balance between native trial storage, a conditional periodic carrier and a reversible receiver.~\cite{baseline}
Its homogeneous native feedback matches carrier work, while receiver loading gives an equal and opposite change of receiver storage. The predecessor proved global exchange evolution, initial supply and reception from empty external stores, and native credit for other prepared carrier states.

Those results leave a distinct dynamical question: does a solution that starts with an empty carrier and receiver later credit the native store without resetting any state? A prepared negative-power state and a reversing symmetry do not answer that question for the original launch.
Here we give an explicit sufficient receiver-coupling threshold for a declared coherent family. The total native power first has the supplying sign and later has the returning sign on the same unique solution. The threshold and witness are finite; propagation and native feedback remain active in the exact equations.

The proof develops the accepted same-trajectory investigation and its internal analytical review.~\cite{investigation,review}
It uses the predecessor's global-solution and launch theorems as premises, then controls a simple comparison by an energy estimate. Section~\ref{sec:problem} specifies the unchanged candidate, Section~\ref{sec:theorem} states the result, and Sections~\ref{sec:energy}--\ref{sec:comparison} give the complete new argument.
The resulting credit is recovery of earlier supplied trial storage. Section~\ref{sec:meaning} distinguishes that fact from a recovered fraction, an attainable capacity or a spatial circuit.

\section{The fixed exchange problem}\label{sec:problem}
\subsection{Chiral drive and finite distributed ports}\label{sec:ports}
Write $x=(a,b)\in\R^6$ and $\Omega=a+ib\in\C^3$. Retain the registration of the predecessor~\cite[Section 1]{baseline}:
\begin{equation}\label{eq:registration}
u=(0,1,0),\quad e_z=(0,0,1),\quad
n_\pm=\frac{2u\pm\sqrt3e_z}{\sqrt7},\quad
A_{\rm reg}=\lambda_gR_z(\delta_g),
\end{equation}
\begin{equation}\label{eq:drive}
K=\kappa_{\rm in}\begin{pmatrix}n_+^T\\n_-^T\end{pmatrix}A_{\rm reg},
\qquad f(x)=K(a\times b),\qquad J(x)=D_xf(x).
\end{equation}
The gains $\lambda_g,\kappa_{\rm in}>0$ and yaw $\delta_g$ are fixed. The real matrix $K$ is $2\times3$ and $J$ is $2\times6$. This is raw chiral projection in the original two-port order. The contact-tangent orientation signs do not insert extra signs in~\eqref{eq:drive}.

Let $s\in\R/L_c\mathbb Z$ on the fixed admissible carrier, with finite $L_c>0$ and speed $c>0$. Its contacts are $s_\pm=L_c/2\mp\eta_{\rm geom}$, with the inherited shortest-separation condition
$\operatorname{dist}_{L_c}(s_+,s_-)=2\eta_{\rm geom}$.
Together with $0<\epsilon_c<\eta_{\rm geom}$, this condition retains disjointness around the whole periodic circle, including across its coordinate seam. Set
\begin{equation}\label{eq:tents}
b_\sigma(s)=\epsilon_c^{-1}
\left(1-\frac{|\operatorname{wrap}_{L_c}(s-s_\sigma)|}{\epsilon_c}\right)_+,
\qquad \sigma\in\{+,-\},
\end{equation}
\begin{equation}\label{eq:ports}
\B v=b_+v_++b_-v_-,\qquad
y=\B^*p=\begin{pmatrix}\int b_+p\dd s\\\int b_-p\dd s\end{pmatrix}.
\end{equation}
All spatial integrals cover one period. These are normalized finite-width tents in $H^1_{\rm per}$. Direct integration of their affine halves and disjoint support give
\begin{equation}\label{eq:profilenorms}
\begin{gathered}
\int b_\sigma\dd s=1,\qquad m=\frac2{3\epsilon_c},\qquad
d=\frac2{\epsilon_c^3},\qquad \B^*\B=mI_2,\\
\norm{\B}_{\R^2\to L^2}=\norm{\B^*}_{L^2\to\R^2}=\sqrt m,\qquad
\norm{\B'v}_2=\sqrt d\,|v|.
\end{gathered}
\end{equation}
The first weak derivative is piecewise constant in $L^2$: continuity of each tent precludes a delta term at its corners. No second derivative or shrinking-width limit is used.

\subsection{State equations, stores and inherited regularity}\label{sec:equations}
Fix $\alpha_N,\beta_R,\eta_N>0$, and write the positive receiver coupling as $\kappa=\eta_R$. Both couplings are dimensionless. For each selected finite $\kappa$, retain
\begin{equation}\label{eq:system}
\boxed{\begin{aligned}
\dot x&=-\frac{\eta_N}{2\alpha_N}J(x)^Ty,\\
p_t+c(r_c)_s&=\B\{\eta_N f(x)-\kappa z\},\\
(r_c)_t+cp_s&=0,\\
\dot z&=\frac{\kappa}{\beta_R}y,\qquad y=\B^*p.
\end{aligned}}
\end{equation}
Here $z\in\R^2$ is one reversible receiver state. All parameters are constant along a solution; both interactions use the same two finite profile regions. We study exchange time only, without native $F$-jumps or a switching protocol.

The trial stores and signed total powers are
\begin{equation}\label{eq:storages}
H_N=\frac{\alpha_N}{2}|x|^2,\qquad
H_C=\frac12\int(p^2+r_c^2)\dd s,\qquad
H_R=\frac{\beta_R}{2}|z|^2,\qquad H_{\rm tot}=H_N+H_C+H_R,
\end{equation}
\begin{equation}\label{eq:powers}
P_{NC}=\eta_N f(x)^Ty,\qquad P_{CR}=\kappa z^Ty.
\end{equation}
Thus $P_{NC}$ is the complete signed sum over the two native ports. A negative individual port does not by itself establish total credit. The positive quadratic $H_N$ differs from the native quartic potential $V$.

The homogeneous identity $J(x)x=2f(x)$, periodic integration and adjointness give the predecessor's accounts~\cite[Section 2]{baseline}:
\begin{equation}\label{eq:balances}
\dot H_N=-P_{NC},\qquad
\dot H_C=P_{NC}-P_{CR},\qquad
\dot H_R=P_{CR},\qquad \dot H_{\rm tot}=0.
\end{equation}
For every fixed finite allowed parameter choice, its Theorem 3.1 supplies a unique global mild solution on
$\R^6\times L^2_{\rm per}{}^2\times\R^2$, with these integrated accounts. If the carrier preparation is in $H^1_{\rm per}{}^2$, then
\begin{equation}\label{eq:regularity}
(p,r_c)\in C(\R;H^1_{\rm per})^2\cap C^1(\R;L^2)^2,\qquad x,z\in C^1(\R).
\end{equation}
We use this established solution theorem. The new energy estimate below is applied to an already existing strong solution; it is not an existence argument based on a formal conservation law.

\subsection{An admitted nonzero coherent launch}\label{sec:launch}
Fix one $0<\phi<\pi/12$ and prepare
\begin{equation}\label{eq:preparation}
\Omega_0=(1,e^{i\phi},1),\qquad p(0)=r_c(0)=z(0)=0.
\end{equation}
The empty carrier is in $H^1$. All moduli are one and the ordinary phase span is $\phi$, as in the predecessor's Proposition 4.1. With $G=\kappa_{\rm in}\lambda_g$,
\begin{equation}\label{eq:gram}
KK^T=G^2\begin{pmatrix}1&1/7\\1/7&1\end{pmatrix},
\qquad R=|x_0|=\sqrt3,\qquad k=\norm K=G\sqrt{8/7}.
\end{equation}
The eigenvalues of the displayed Gram matrix are $8G^2/7$ and $6G^2/7$. Also $a_0\times b_0=(-\sin\phi,0,\sin\phi)$, so
\begin{equation}\label{eq:launchdrive}
f_{0,\pm}=\frac{G\sin\phi}{\sqrt7}(-2\sin\delta_g\pm\sqrt3),
\qquad |f_0|^2=\frac{G^2\sin^2\phi}{7}(6+8\sin^2\delta_g)>0.
\end{equation}
The two entries differ by $2\sqrt3G\sin\phi/\sqrt7\ne0$ for every fixed yaw. An individual entry may vanish, but the vector does not. No invariance of this launch family or its coherent phase domain is assumed.

\section{A sufficient finite-coupling theorem}\label{sec:theorem}
Fix the admissible geometry and finite profiles of Section~\ref{sec:ports}, the speed, registration, positive gains and storage scales, a positive $\eta_N$, and one phase $\phi$ in~\eqref{eq:preparation}. These choices precede the choice of receiver coupling. In addition to $R,k,m,d$ above, define
\begin{equation}\label{eq:constants}
F_{\max}=\frac{kR^2}{2},\qquad
\omega=\sqrt{\frac m{\beta_R}},\qquad
T=\frac{3\pi}{2\omega},\qquad
D=\frac m\omega|f_0|^2,
\end{equation}
\begin{equation}\label{eq:errorconstants}
\begin{aligned}
A&=\frac{\eta_N^2k^2R^2mF_{\max}}{4\alpha_N},\\
C_1&=\frac{3c\sqrt{md}\,|f_0|^2}{\omega^2},&
C_2&=mAT^3\left(F_{\max}+\frac{|f_0|}{3}\right),
\end{aligned}
\end{equation}
\begin{equation}\label{eq:threshold}
\boxed{\kappa_{\rm crit}
=\max\left\{1,\frac{4C_1}{D},\sqrt{\frac{4C_2}{D}}\right\}.}
\end{equation}
Every quantity in~\eqref{eq:constants}--\eqref{eq:threshold} is independent of $\kappa$. They are finite, and $D>0$ by~\eqref{eq:launchdrive}. Appendix~\ref{sec:units} checks their dimensions.

\Needspace{7\baselineskip}
\begin{theorem}[Same-trajectory native credit]\label{thm:credit}
For each fixed finite $\kappa\ge\kappa_{\rm crit}$, the unique global solution of~\eqref{eq:system} with preparation~\eqref{eq:preparation} has the following properties.

For all sufficiently small positive time, its total native power, receiver storage and receiver power satisfy
$P_{NC}(t)>0$, $H_R(t)>0$ and $P_{CR}(t)>0$.
At the explicit later time
\begin{equation}\label{eq:witness}
t_*=\frac T\kappa,\qquad
P_{NC}(t_*)\le-\frac{\eta_N^2D}{2\kappa}<0.
\end{equation}
There exists $0<\rho_\kappa<t_*/4$ such that
\begin{equation}\label{eq:interval}
P_{NC}(t)\le-\frac{\eta_N^2D}{4\kappa}
\qquad(t\in[t_*-\rho_\kappa,t_*+\rho_\kappa]),
\end{equation}
and consequently
\begin{equation}\label{eq:gain}
H_N(t_*+\rho_\kappa)-H_N(t_*-\rho_\kappa)
\ge\frac{\eta_N^2D\rho_\kappa}{2\kappa}>0.
\end{equation}
The initial positive interval can be shortened to lie strictly before this later interval. Both occur on the same solution with constant parameters and no reset.
\end{theorem}

The threshold is an explicit sufficient condition, not a necessary or optimal one. The witness and power margin are quantitative; the interval width is supplied by continuity. Neither a uniform lower bound for $\rho_\kappa$ nor a substantial recovered amount follows from the theorem. Its proof occupies the next two sections.

\section{Scaled energy and native drift}\label{sec:energy}
\subsection{Exact finite-coupling scaling}\label{sec:scaling}
For the selected finite $\kappa>0$, put
\begin{equation}\label{eq:scaling}
\tau=\kappa t,\quad x_\kappa(\tau)=x(\tau/\kappa),\quad
(p,r_c,z)(\tau/\kappa)=\frac{\eta_N}{\kappa}
(\widehat p,\widehat r,\widehat z)(\tau),\quad
\widehat y=\B^*\widehat p.
\end{equation}
The chain rule gives the exact system
\begin{equation}\label{eq:scaledsystem}
\begin{aligned}
(x_\kappa)_\tau&=-\frac{\eta_N^2}{2\alpha_N\kappa^2}J(x_\kappa)^T\widehat y,\\
\widehat p_\tau+\frac c\kappa\widehat r_s&=\B\{f(x_\kappa)-\widehat z\},\\
\widehat r_\tau+\frac c\kappa\widehat p_s&=0,\\
\widehat z_\tau&=\frac1{\beta_R}\widehat y,
\end{aligned}
\end{equation}
and the original total power is
\begin{equation}\label{eq:powerconversion}
P_{NC}(\tau/\kappa)=\frac{\eta_N^2}{\kappa}
f(x_\kappa(\tau))^T\widehat y(\tau).
\end{equation}
The scaled time still has time units because $\kappa$ is dimensionless. This change of variables leaves the original propagation speed $c>0$ unchanged.

\subsection{The weighted external estimate}\label{sec:forced}
On $\EE=L^2_{\rm per}\times L^2_{\rm per}\times\R^2$, set
\begin{equation}\label{eq:externalnorm}
\norm{W}_{CR}^2=\norm q_2^2+\norm h_2^2+\beta_R|v|^2,
\qquad W=(q,h,v),
\end{equation}
\begin{equation}\label{eq:operator}
\begin{aligned}
\LL_\kappa W&=
\left(-\frac c\kappa h_s-\B v,\,
      -\frac c\kappa q_s,\,
      \frac1{\beta_R}\B^*q\right),\\
D(\LL_\kappa)&=H^1_{\rm per}{}^2\times\R^2.
\end{aligned}
\end{equation}
For a domain element, periodic integration by parts and adjointness give
\begin{equation}\label{eq:skew}
\begin{aligned}
\langle W,\LL_\kappa W\rangle_{CR}
&=-\frac c\kappa\int(qh_s+hq_s)\dd s\\
&\phantom{={}}-\langle q,\B v\rangle+v^T\B^*q=0.
\end{aligned}
\end{equation}
Thus, for a strong curve $W_\tau=\LL_\kappa W+g$,
\begin{equation}\label{eq:forcedbound}
\norm{W(\tau)}_{CR}\le\norm{W(0)}_{CR}
+\int_0^\tau\norm{g(s)}_{CR}\dd s,\qquad \tau\ge0.
\end{equation}
At a zero of the norm, differentiate
$\sqrt{\norm W_{CR}^2+\varepsilon^2}$ instead; its derivative is at most $\norm g_{CR}$ by~\eqref{eq:skew}. Integrating and letting $\varepsilon\downarrow0$ proves~\eqref{eq:forcedbound} without division by zero. The homogeneous linear part preserves this norm. We use the estimate only for the known strong solution and its difference from the explicit comparison, so no additional group-limit assertion is needed.

Since the external stores start empty, conservation~\eqref{eq:balances} implies $|x_\kappa(\tau)|\le R$. The derivative formula
$J(x)h=K(h_a\times b+a\times h_b)$ gives
\begin{equation}\label{eq:ballbounds}
|f(x)|\le\frac k2|x|^2,\qquad
\norm{J(x)}\le k|x|,\qquad
|f(x)-f(\widetilde x)|\le kR|x-\widetilde x|
\end{equation}
when $x,\widetilde x$ lie in the ball of radius $R$. The last inequality follows by integrating $J$ along the line segment in that convex ball. No preservation of a coherent phase domain enters.

Apply~\eqref{eq:forcedbound} to
$V_\kappa=(\widehat p,\widehat r,\widehat z)$ with forcing
$(\B f(x_\kappa),0,0)$ and zero initial value. Equations~\eqref{eq:profilenorms} and~\eqref{eq:ballbounds} yield, for every finite $\tau\ge0$,
\begin{equation}\label{eq:externalbounds}
\norm{V_\kappa(\tau)}_{CR}\le\sqrt mF_{\max}\tau,\qquad
|\widehat y(\tau)|\le\sqrt m\norm{\widehat p(\tau)}_2
\le mF_{\max}\tau.
\end{equation}
These constants do not depend on $\kappa$. Retaining receiver loading inside the norm-preserving linear part is essential to the estimate.

\subsection{The native drift bound}\label{sec:drift}
Integrating the first equation in~\eqref{eq:scaledsystem}, with~\eqref{eq:ballbounds}--\eqref{eq:externalbounds}, gives
\begin{equation}\label{eq:coredrift}
|x_\kappa(\tau)-x_0|
\le\frac{\eta_N^2kR}{2\alpha_N\kappa^2}
\int_0^\tau mF_{\max}s\dd s
=\frac{\eta_N^2kRmF_{\max}}{4\alpha_N\kappa^2}\tau^2.
\end{equation}
The convex-ball Lipschitz bound consequently implies
\begin{equation}\label{eq:drivedrift}
|f(x_\kappa(\tau))-f_0|\le\frac{A\tau^2}{\kappa^2}.
\end{equation}
The factor one-fourth includes the native feedback factor one-half and the integral of time. These are bounds for the exact finite-coupling solution, established before introducing the comparison.

\section{Comparison, finite error and the returning sign}\label{sec:comparison}
\begin{proof}[Proof of Theorem~\ref{thm:credit}]
\textit{The explicit comparison.}
Freeze the core and omit scaled propagation only in the comparison problem. Its zero-external-state solution is
\begin{equation}\label{eq:comparison}
\begin{aligned}
x^{(0)}(\tau)&=x_0,&
\widehat p^{(0)}(\tau)&=\frac{\sin(\omega\tau)}{\omega}\B f_0,\\
\widehat r^{(0)}(\tau)&=0,&
\widehat z^{(0)}(\tau)&=f_0\{1-\cos(\omega\tau)\},
\end{aligned}
\end{equation}
\begin{equation}\label{eq:comparisonoutput}
\widehat y^{(0)}(\tau)=\frac m\omega f_0\sin(\omega\tau).
\end{equation}
Indeed $\widehat p^{(0)}_\tau=\B(f_0-\widehat z^{(0)})$ and
$\widehat z^{(0)}_\tau=\widehat y^{(0)}/\beta_R$, because
$\omega^2=m/\beta_R$. Every external coordinate vanishes initially. At the specified $T$,
\begin{equation}\label{eq:comparisonmargin}
Q_0(T):=f_0^T\widehat y^{(0)}(T)
=-\frac m\omega|f_0|^2=-D<0.
\end{equation}
This is a comparison margin; finite error control is required to transfer its sign to the actual solution.

\textit{Error equation and regularity.}
Let $V_0=(\widehat p^{(0)},0,\widehat z^{(0)})$ and
$E_\kappa=V_\kappa-V_0$. The exact solution has regularity~\eqref{eq:regularity}, preserved by the finite scaling. The explicit comparison is also a strong curve, with carrier in $C^1_\tau H^1_s$. Hence their difference already belongs to $C_\tau H^1_s\cap C^1_\tau L^2_s$ in its carrier coordinates. This justifies its energy calculation before an estimate from the residual is used.

Subtract the comparison equations while retaining the finite operator~\eqref{eq:operator} on the error:
\begin{equation}\label{eq:error}
\begin{aligned}
(E_\kappa)_\tau={}&\LL_\kappa E_\kappa
+(\B(f(x_\kappa)-f_0),0,0)\\
&+\left(0,-\frac c\kappa\B'f_0
                 \frac{\sin(\omega\tau)}{\omega},0\right),
\qquad E_\kappa(0)=0.
\end{aligned}
\end{equation}
The propagation residual is negative in the second carrier component:
it is $-(c/\kappa)\partial_s\widehat p^{(0)}$.
The first component has no propagation residual because
$\widehat r^{(0)}=0$. By~\eqref{eq:profilenorms}, the residual lies in $L^2$ and involves only the first weak derivative of a tent. It need not belong to $H^1$; strong regularity of the error follows from subtraction, not from differentiating that residual. No second tent derivative, Fourier truncation or operator-norm convergence of an unbounded generator is invoked.

Apply~\eqref{eq:forcedbound} to~\eqref{eq:error}, using~\eqref{eq:drivedrift}:
\begin{equation}\label{eq:errorintegral}
\norm{E_\kappa(\tau)}_{CR}
\le\frac{c\sqrt d\,|f_0|}{\kappa\omega}
\int_0^\tau|\sin(\omega s)|\dd s
+\frac{\sqrt m A\tau^3}{3\kappa^2},
\qquad 0\le\tau\le T.
\end{equation}
Splitting the absolute sine integral at $\pi/\omega$ gives
$2/\omega+1/\omega=3/\omega$ at $T$. Therefore
\begin{equation}\label{eq:finiteerror}
\boxed{\norm{E_\kappa(T)}_{CR}
\le\frac{3c\sqrt d\,|f_0|}{\kappa\omega^2}
+\frac{\sqrt m AT^3}{3\kappa^2}.}
\end{equation}

\textit{The full two-port power error.}
Write $Q_\kappa=f(x_\kappa(T))^T\widehat y(T)$. The exact decomposition is
\begin{equation}\label{eq:powerdecomposition}
\begin{aligned}
Q_\kappa-Q_0(T)
={}&(f(x_\kappa(T))-f_0)^T\widehat y(T)\\
&+f_0^T\{\widehat y(T)-\widehat y^{(0)}(T)\}.
\end{aligned}
\end{equation}
The first term uses the actual output. The bounded adjoint,~\eqref{eq:externalbounds},~\eqref{eq:drivedrift} and~\eqref{eq:finiteerror} give
\begin{equation}\label{eq:powererror}
\begin{aligned}
|Q_\kappa+D|
&\le\frac{AT^2}{\kappa^2}\,mF_{\max}T\\
&\quad+|f_0|\sqrt m\left\{
\frac{3c\sqrt d\,|f_0|}{\kappa\omega^2}
+\frac{\sqrt m AT^3}{3\kappa^2}\right\}\\
&=\frac{C_1}{\kappa}+\frac{C_2}{\kappa^2}.
\end{aligned}
\end{equation}
This controls the complete signed sum, including both product-error contributions. For $\kappa\ge\kappa_{\rm crit}$,
\begin{equation}\label{eq:allocation}
\frac{C_1}{\kappa}\le\frac D4,\qquad
\frac{C_2}{\kappa^2}\le\frac D4,\qquad
Q_\kappa\le-\frac D2.
\end{equation}
Converting by~\eqref{eq:powerconversion} proves the finite original-time witness~\eqref{eq:witness}.

\textit{Continuity, initial activity and chronology.}
The native power is continuous: $x$ is continuous, $f$ is polynomial, $p$ is continuous in $L^2$, and $\B^*$ is bounded. Its strict margin at $t_*$ permits
$0<\rho_\kappa<t_*/4$ with~\eqref{eq:interval} on the whole closed interval. Integrating $\dot H_N=-P_{NC}$ over its length $2\rho_\kappa$ proves~\eqref{eq:gain}.

For this same fixed finite $\kappa$, the predecessor's Proposition 4.1 applies because $\eta_N,\kappa>0$ and $f_0\ne0$. Its empty-launch expansions give
\begin{equation}\label{eq:initialactivity}
\begin{aligned}
P_{NC}(t)&=\eta_N^2m|f_0|^2t+o(t),\\
H_R(t)&=\frac{\eta_N^2\kappa^2m^2|f_0|^2}{8\beta_R}t^4+o(t^4),\\
P_{CR}(t)&=\frac{\eta_N^2\kappa^2m^2|f_0|^2}{2\beta_R}t^3+o(t^3).
\end{aligned}
\end{equation}
All leading coefficients are positive. Shorten their common positive interval to lie in $(0,t_*/4)$; the later interval lies in
$(3t_*/4,5t_*/4)$. The remainders need not be uniform in $\kappa$, and we never evaluate these short-time expansions at $T/\kappa$. Global uniqueness identifies the initial activity and later credit as stages of one unchanged solution. Neither a prepared negative carrier state nor a reversing involution is used.
\end{proof}

\section{Recovery, scale and invariant capacity}\label{sec:meaning}
\subsection{Recovery without cumulative native gain}\label{sec:account}
The returning sign in Theorem~\ref{thm:credit} has an exact storage meaning. Empty external preparation and~\eqref{eq:balances} give, for all future exchange times,
\begin{equation}\label{eq:cumulative}
H_N(0)-H_N(t)=H_C(t)+H_R(t)
=\int_0^tP_{NC}(s)\dd s\ge0.
\end{equation}
Thus $H_N(t)\le H_N(0)$ even during a later credit interval.
The positive increment~\eqref{eq:gain} recovers earlier supplied trial storage; it does not create cumulative native gain above the initial stock.

At the comparison witness,
\begin{equation}\label{eq:comparisonbackflow}
\widehat z^{(0)}(T)=f_0,\qquad
\widehat y^{(0)}(T)=-\frac m\omega f_0.
\end{equation}
Their negative dot product identifies receiver backflow in that comparison. The finite threshold proves exact total native credit and initial exact receiver activity. It does not separately prove $P_{CR}(t_*)<0$ for every coupling satisfying~\eqref{eq:threshold}; an assertion about that different power requires its own finite bound.

\subsection{Fixed scaled windows and transferred amounts}\label{sec:scales}
For any fixed finite $S>0$, scaling and~\eqref{eq:externalbounds} imply
\begin{equation}\label{eq:stockscale}
\begin{aligned}
H_C(\tau/\kappa)+H_R(\tau/\kappa)
&=\frac{\eta_N^2}{2\kappa^2}\norm{V_\kappa(\tau)}_{CR}^2\\
&\le\frac{\eta_N^2mF_{\max}^2S^2}{2\kappa^2},
\qquad 0\le\tau\le S.
\end{aligned}
\end{equation}
Taking $S=T$ covers launch through the witness; $S=5T/4$ also covers the later interval. With all other parameters fixed, the external stock is
$O(\eta_N^2/\kappa^2)$ on such scaled windows. Equation~\eqref{eq:powererror} similarly gives
\begin{equation}\label{eq:powerscale}
P_{NC}(t_*)=-\frac{\eta_N^2D}{\kappa}
+O\!\left(\frac{\eta_N^2}{\kappa^2}\right).
\end{equation}
These are large-coupling statements on fixed scaled windows, not long-time decay statements. They show that strict existence of credit can coexist with a small absolute transferred stock. A small stock does not determine the fraction of it recovered: that question requires a specified numerator, denominator and time window. The present theorem establishes neither a substantial returned share nor maximum accessible extraction.

\subsection{The unchanged invariant ceiling}\label{sec:capacity}
The predecessor's Theorem 5.1 preserves the complex bilinear quantity
\begin{equation}\label{eq:invariant}
Z_2(\Omega)=\sum_{i=1}^3\Omega_i^2,\qquad
I=|\Omega|^2,\qquad
H_N(t)\ge\frac{\alpha_N}{2}|Z_2(\Omega_0)|.
\end{equation}
For the initial phase span $\phi$, combining its net-extraction bound with~\eqref{eq:cumulative} yields
\begin{equation}\label{eq:capacity}
\begin{aligned}
0\le\frac{H_N(0)-H_N(t)}{H_N(0)}
&\le\frac{I_0-|Z_2(\Omega_0)|}{I_0}
\le1-\cos\phi\\
&<1-\cos(\pi/12)\simeq0.0340741737109317.
\end{aligned}
\end{equation}
The denominator is positive initial core trial storage. The approximately
$3.4074\%$ maximum-span ceiling is a bound on net extraction, not an achieved fraction, gross repeated work, receiver share or diversion percentage. Its persistence needs invariance of $Z_2$, not continued membership in the initial coherent domain. A later native gain is entirely compatible with this ceiling.

\section{Geometric meaning and conclusion}\label{sec:conclusion}
The explicit witness satisfies
\begin{equation}\label{eq:traveltime}
\frac{t_*}{L_c/c}=\frac{cT}{\kappa L_c}\longrightarrow0
\qquad(\kappa\to\infty),
\end{equation}
with geometry and speed fixed. Since both couplings occupy the same finite profiles, this certificate describes fast collocated exchange. It does not show that a wave has completed a circuit before credit occurs.

The motivating picture is a central starting/meeting region, opposed cones, outward spreading and curvature around an intended tight toroidal arrangement, inward return and possible transfer elsewhere. Stationary supports remain stationary. The adopted one-dimensional circle, intended tight toroid, close original sphere and extended-line sphere are distinct; the finite host-notch $N$ regions differ from the six central broken $G$ boundaries.~\cite[Section 8]{baseline} Exclusion of exact $O$ is not exclusion of every finite central region. The theorem selects no new spatial receiver, six-opening map or field interpretation. The author's tentative \emph{15--20\%, possibly more} remains geometric motivation rather than a fitted target or a consequence of the initial-core extraction bound.

Within the unchanged reciprocal chiral candidate, the dynamical gap between initial empty-launch supply and separately prepared credit is now closed at an explicit sufficient finite-coupling scope. The exact receiver is active at launch, all coefficients remain constant, and the total native power later becomes negative on that same solution. The proof controls propagation and core drift with finite constants and first weak derivatives; it does not replace the exact trajectory by its comparison. Recovery remains below the initial native stock, and neither arbitrary-coupling classification nor a substantial recovered fraction or physical spatial circuit is established.

The next scientific questions concern accessible transfer and receiver behavior, followed by native stepping/cadence, intended route correspondence and six-opening transport. Gravity--magnetism remains linked across that programme; outer Tri-Octagon geometry and further justified SRG integration remain later directions. This manuscript and its review/publication stage precede those investigations.

\subsection*{Contributions and evidence}
Hilmir Fr\'imann Halld\'orsson (pzychozen) supplies research authorship and direction. GPT formulated and internally accepted the underlying finite-coupling investigation; Codex independently verified it and prepared this companion manuscript.~\cite{investigation,review}
The supporting record contains 25 distinct exact-algebra checks, independently replayed, not 50 distinct identities. These support algebra and coefficient collection; the written argument and cited foundations justify regularity, inequalities, admissibility and continuity. No scientific replay or trajectory computation was performed for manuscript preparation. Internal acceptance of the investigation is distinct from review of this final typeset manuscript, which is pending, and from external peer review or physical validation.

\appendix
\section{Dimensions of the certificate}\label{sec:units}
Use the predecessor's bookkeeping native amplitude unit $\mathrm A$ and abstract carrier units $\mathrm U,\mathrm L,\mathrm T$. The normalized launch uses unit native amplitudes. The drift constant $A$ in~\eqref{eq:errorconstants} is distinct from both the unit $\mathrm A$ and the registration matrix $A_{\rm reg}$. The inherited and derived dimensions are
\begin{center}\small
\begin{tabularx}{\textwidth}{@{}P{.45\textwidth}Y@{}}\toprule
Quantity & Dimension\\\midrule
$x,R;\ K,k;\ J$ & $\mathrm A;\ \mathrm U\mathrm L/(\mathrm T\mathrm A^2);\ \mathrm U\mathrm L/(\mathrm T\mathrm A)$\\[3pt]
$p,r_c,y;\ f,z,F_{\max}$ & $\mathrm U;\ \mathrm U\mathrm L/\mathrm T$\\[3pt]
$m;\ d;\ c$ & $\mathrm L^{-1};\ \mathrm L^{-3};\ \mathrm L/\mathrm T$\\[3pt]
$\alpha_N;\ \beta_R;\ \omega$ & $\mathrm U^2\mathrm L/\mathrm A^2;\ \mathrm T^2/\mathrm L;\ \mathrm T^{-1}$\\[3pt]
$\tau,T,t_*,\rho_\kappa;\ \eta_N,\kappa$ & $\mathrm T;\ 1$\\[3pt]
$A;\ \norm{V}_{CR}$ & $\mathrm U\mathrm L/\mathrm T^3;\ \mathrm U\mathrm L^{1/2}$\\[3pt]
$H_N,H_C,H_R$ & $\mathrm U^2\mathrm L$\\[3pt]
$D,C_1,C_2,Q_\kappa,P_{NC},P_{CR}$ & $\mathrm U^2\mathrm L/\mathrm T$\\[3pt]
$C_1/D,\ C_2/D,\ \kappa_{\rm crit}$ & $1$\\\bottomrule
\end{tabularx}
\end{center}
In particular the terms $c\sqrt d|f_0|/\omega^2$ and
$\sqrt m AT^3$ in~\eqref{eq:finiteerror} both have external-norm units. Both threshold ratios are dimensionless, as required for a dimensionless coupling. These checks establish mathematical commensurability, not physical calibration.

\begin{thebibliography}{9}
\bibitem{baseline}
H.~F.~Halld\'orsson,
\emph{A Reciprocal Chiral Exchange Extension of the Tri-Octagon Model:
Well-posedness, matched storage, and extraction limits}.
Author-authorized repository research edition v0.1, 3 November 2026.
\href{https://github.com/pzychozen/trioctagon-physics/tree/273b1a6faaab19ea58748d00576439166eeef9ab/papers/RECIPROCAL_CHIRAL_EXCHANGE/publication/v0.1}{Immutable edition and editable sources}.
The present paper uses Section~1 (candidate and units), Section~2 (balances),
Theorem~3.1 (global solutions), Proposition~4.1 (empty launch),
Theorem~5.1 (invariant capacity) and Section~8 (interpretation).

\bibitem{investigation}
H.~F.~Halld\'orsson (research direction), with independent analytical preparation by Codex,
\emph{Same-trajectory native credit from an initially empty launch}.
Technical investigation NATIVE\_SAME\_TRAJECTORY\_RETURN\_01,
3 November 2026, equations R1--R30 and ledger STR01--STR28.
\href{evidence/investigation/NATIVE_SAME_TRAJECTORY_RETURN_01_PROOF.md}{Supporting proof report};
\href{evidence/investigation/NATIVE_SAME_TRAJECTORY_RETURN_01_CLAIM_SOURCE_LEDGER.md}{original claim/source ledger}.
Unchanged technical evidence accompanying this manuscript.

\bibitem{review}
GPT,
\emph{Review of same-trajectory native return}.
Project-internal mathematical and evidence review,
3 November 2026, Sections~1--5.
\href{evidence/current_gpt_review.md}{Supporting analytical acceptance}.
This record accepts the underlying investigation and unchanged constants;
it is not a review of the present final PDF.
\end{thebibliography}

\end{document}
